\begin{abstract}
Large language models are increasingly becoming intermediaries in healthcare communication. However, most existing evaluations focus on the quality of generated answers and implicitly assume that the user's underlying concern has already been correctly understood. This assumption is particularly problematic in women's-health communication, where the same concern may be expressed through clinical terminology, everyday language, or indirect descriptions (e.g., ``irregular menstrual bleeding,'' ``my period is all over the place,'' or ``something about my cycle feels wrong''), while equivalent concerns may be communicated differently across languages and cultures. To address this gap, we introduce \textbf{HerHealthEval}, a multilingual framework for evaluating language understanding in women's-health communication. HerHealthEval extracts core clinical concerns from existing cases and generates matched variants across five registers (clinical, lay, indirect, ambiguous, and emotional), localized into English, French, and Modern Standard Arabic through native-speaker linguistic validation. The resulting cases are split into adaptation and evaluation sets using a strict leakage-control protocol that prevents source and paraphrase-family overlap. %Evaluating a baseline multilingual model against multilingual domain-adapted variants reveals critical safety failures hidden by standard metrics, including clinical under-triage, a near-total collapse of clarification-seeking behavior, and severe language-specific regressions.
Comparisons between a baseline multilingual model and domain-adapted variants reveal critical safety failures hidden by standard metrics, including clinical under-triage, collapse of clarification-seeking behavior, and severe language-specific regressions. In particular, multilingual adaptation improves apparent cross-language consistency while simultaneously suppressing clarification-seeking behavior to nearly zero and increasing under-triage rates in both French and Arabic to 0.994. We trace this collapse to a language-asymmetric risk-labeling defect rather than to multilingual adaptation itself, and show that re-adapting with risk labels derived once from the English source restores non-English under-triage to below the baseline (0.994$\rightarrow$0.572/0.558), while clarification-seeking remains suppressed. Collectively, these findings suggest that multilingual healthcare systems require evaluating language understanding separately from response generation and motivate guardrailed, uncertainty-aware dialogue systems for women's-health communication.
\end{abstract}

